\documentclass[12pt]{article}
\usepackage{amsmath, amssymb, booktabs, geometry}
\geometry{margin=1in}
\setlength{\parindent}{0pt}
\setlength{\parskip}{0.8em}

\title{ECON 101: Macroeconomics \\ Quiz 3}
\author{Instructor: Muhebullah Karimzada}
\date{October 22,  2025}

\begin{document}
\maketitle
\hrule
\vspace{1em}

\section*{Multiple Choice Questions (1 mark each)}

\textbf{1.}  
The Consumer Price Index (CPI) in 2022 was 152 and in 2023 it rose to 160.  
What was the annual inflation rate?  

(A) 4.8\% \quad (B) 5.3\% \quad (C) 6.0\% \quad (D) 8.0\%


\textbf{2.}  
If the nominal interest rate is 5\% and the inflation rate is 3\%, the real interest rate is:  

(A) 8\% \quad (B) 2\% \quad (C) 1.5\% \quad (D) –2\%



\textbf{3.}  
Which of the following items is \emph{not} included in the Consumer Price Index basket?  

(A) Imported televisions \quad (B) Used cars \quad (C) Restaurant meals \quad (D) Electricity bills



\textbf{4.}  
Suppose real GDP per capita in Country A is \$42{,}000 and in Country B is \$28{,}000.  
If Country A grows at 2\% per year while B grows at 4\% per year, approximately how many years will it take for B’s income to catch up with A’s?  

(A) 21 years \quad (B) 26 years \quad (C) 35 years \quad (D) 42 years



\textbf{5.}  
During a business-cycle recession, which variable usually falls first?  

(A) Unemployment \quad (B) Inflation \quad (C) Interest rates \quad (D) Investment spending


\newpage
\section*{Problem (15 marks)}

\textbf{6. Economic Growth and Inflation Analysis}

The following data describe a small economy over two years:

\begin{center}
\begin{tabular}{lcccc}
\toprule
\textbf{Item} & \textbf{Quantity (Year 1)} & \textbf{Quantity (Year 2)} & \textbf{Price (Year 1)} & \textbf{Price (Year 2)} \\
\midrule
Smartphones & 200 & 250 & \$500 & \$520 \\
Tablets & 100 & 120 & \$800 & \$880 \\
Headphones & 300 & 330 & \$100 & \$90 \\
\bottomrule
\end{tabular}
\end{center}

\textbf{Tasks:}
\begin{enumerate}
    \item[(a)] Using \textbf{Year 1 as the base year}, calculate:
        \begin{enumerate}
            \item Nominal GDP for Year 1 and Year 2.
            \item Real GDP for Year 2 in Year 1 prices.
            \item The real GDP growth rate.
        \end{enumerate}

    \item[(b)] Compute the \textbf{GDP deflator} for each year and the corresponding \textbf{inflation rate}.

    \item[(c)] Suppose the population rises from 800 people in Year 1 to 820 people in Year 2.  
    Calculate the growth rate of real GDP per capita.

    \item[(d)] Based on your results, briefly discuss whether living standards improved.
\end{enumerate}


\vspace{1em}
\hrule
\vspace{7 cm}

\begin{center}
\textit{End of Quiz}
\end{center}

\end{document}
